\documentclass[10pt,conference]{IEEEtran}
\usepackage{booktabs}
\usepackage{hyperref}
\usepackage{xurl}

\title{Replayable State-Diff Harnesses for Repository-Changing Agents}
\author{\IEEEauthorblockN{Song Luo}
\IEEEauthorblockA{\texttt{luosongred@gmail.com}}}

\begin{document}
\maketitle

\begin{abstract}
Software agents now edit files, run tests, inspect repositories, and prepare
changes that may affect downstream maintenance. This tool demonstration
presents a replayable harness for evaluating whether such agents operate under
explicit state and permission contracts. The artifact models tool calls,
observations, permission gates, memory scope, replay, and final state diffs. It
contains 24 cases, raw run logs, regenerated tables, and validation scripts. For
software analysis and evolution, the key contribution is an inspectable way to
treat an agent trajectory as a maintainable software artifact: what changed,
which tool produced it, whether policy allowed it, and whether the failure can
be replayed. Tool: \url{https://github.com/rrrrrredy/agent-harness-paper}.
Archive: \url{https://doi.org/10.5281/zenodo.20907471}. Video:
\url{https://github.com/rrrrrredy/agent-harness-paper/releases/download/v0.1.0-preprint/agent-harness-saner-tool-demo-video.mp4}.
\end{abstract}

\section{Problem}

Coding and workflow agents create new maintenance questions. A final answer
does not reveal whether the agent edited the intended file, respected a dry-run
policy, avoided secrets, scoped memory correctly, or left a replayable trace.
These questions are central to software analysis, evolution, and
reengineering because agent behavior becomes part of the evolving development
process.

\section{Tool}

The demonstrated tool is a small replayable harness for stateful agent
contracts. It provides:

\begin{itemize}
\item case definitions with initial state, allowed tools, forbidden tools, and
state assertions;
\item mock tools for file, repository, memory, email/browser-like, and review
state;
\item prompt variants for no harness, thick checklist, and thin contract;
\item metrics for task success, state-diff correctness, invalid tool calls,
permission violations, unsafe secret access, and replayability;
\item scripts for validation, regeneration, source evidence audit, and secret
scanning.
\end{itemize}

\section{Demo}

The demo walks through a repository-changing scenario:

\begin{enumerate}
\item load a case that requires a safe repository update;
\item show the expected tool contract and forbidden actions;
\item run the harness and inspect the resulting action trace;
\item compare initial and final mock state;
\item inspect a failure where a broad prompt checklist does not enforce the
contract;
\item regenerate aggregate tables from raw runs.
\end{enumerate}

The artifact is designed for reviewers to execute without live model keys. Live
provider runs are included as raw rows; deterministic validation checks the
simulator and metrics.

\section{Relevance to SANER}

The harness treats agent trajectories as analyzable software artifacts. It can
support regression testing of agent workflows, analysis of repository side
effects, evolution of prompt and tool schemas, and reengineering of brittle
agent scripts into explicit contracts. This fits SANER's focus on extracting
useful information from software and system artifacts and using that information
for maintenance and renewal.

\section{Availability}

The repository and release archive are public. Code is MIT licensed; manuscript,
benchmark cases, evidence snapshots, tables, figures, and documentation are CC
BY 4.0.

\begin{thebibliography}{9}
\bibitem{toolsandbox} ToolSandbox, stateful tool-use evaluation benchmark.
\bibitem{agentdojo} AgentDojo, dynamic tool-calling and prompt-injection evaluation.
\bibitem{sweagent} SWE-agent and agent-computer interfaces for software engineering agents.
\bibitem{deli} Deli AutoResearch framework for long-horizon autonomous tasks.
\end{thebibliography}

\end{document}
